\documentclass[11pt,twocolumn]{article}

\usepackage[utf8]{inputenc}
\usepackage[T1]{fontenc}
\usepackage{amsmath,amssymb,amsfonts}
\usepackage{booktabs}
\usepackage{hyperref}
\usepackage[margin=0.75in]{geometry}

\hypersetup{colorlinks=true,linkcolor=black,citecolor=black,urlcolor=blue}

\title{\textbf{Large Pattern Models: Replacing Token Consumption\\
with Artifact Reuse}}
\author{iulian dacineu\\
\textit{unia} --- \texttt{https://github.com/dacineu/unia}}
\date{September 2026}

\begin{document}
\maketitle

\begin{abstract}
Language models are consumed per request. A coding agent that has solved a task
before pays for it again, because the solution it produced is stored as
\emph{text} --- in context, in a commit, in a skill file --- and text must be
re-read by a transformer to become an action. This paper argues that the unit of
reuse should be an \emph{executable artifact} rather than a \emph{description},
and that the resulting system should be optimised for a metric we call
\textbf{escalation rate}: the fraction of requests that cannot be served from a
local artifact and must reach a provider.

I formalise a resource model, the \texttt{.ure} format, in which a unit of
capability is a content-addressed declaration bound to a runnable body. I
describe an induction loop that converts recorded interactions into candidate
artifacts without gradient descent, an evidence gate that makes premature
promotion structurally difficult, and a collection protocol in which each
artifact observes itself. The prototype is \emph{unia}, in Rust, with a Model
Context Protocol server surface.

Measurements characterise the retrieval tier and establish a first
retrieval-quality baseline: a $4.6\,\mu$s lexical matcher achieves 90\% Hit@1
against a hand-written fixture, with all observed failures attributable to
manifest authoring rather than to a need for semantic machinery. I am explicit
that this fixture is too small to support the thesis, and I state the experiment
that would falsify it.
\end{abstract}

\section{Introduction}

The dominant cost of a language-model-backed agent is not reasoning. It is
re-deriving knowledge the system already possesses.

Consider a coding agent asked to add a dependency and refresh a lockfile. The
second time it is asked, the answer is identical. The first answer may have
taken 2{,}000 input tokens and 600 output tokens. The second costs the same,
because the first answer was recorded as \emph{text} --- a diff, a file, a
markdown skill --- and text is inert. It becomes an action only by passing
through a transformer again.

Three mechanisms mitigate this, and all three leave the architecture unchanged:

\begin{center}
\small
\begin{tabular}{@{}ll@{}}
\toprule
\textbf{Mechanism} & \textbf{Eliminates} \\
\midrule
Prompt caching & input cost on a repeat \\
Semantic caching & the request, on near-match \\
RAG & nothing; it adds context \\
\bottomrule
\end{tabular}
\normalsize
\end{center}

The ceiling on all three is set by the same thing. A cache can only return
something previously computed and stored. If the stored artifact is a
\emph{description} of an action, serving it requires a model. If the stored
artifact \emph{is} the action, serving it does not.

\subsection{The claim}

\begin{quote}
\textbf{C1.} For a workload with sufficient repetition, a system that stores
executable artifacts rather than descriptions can drive escalation rate toward
zero without gradient training, and per-served-request token cost toward zero,
because the model is \emph{out-routed} rather than \emph{distilled}.
\end{quote}

``Sufficient repetition'' without training is a deliberate constraint, not a
limitation: the learned artifact is inspectable, auditable, and deletable.

\subsection{What would falsify C1}

C1 fails if, on a labelled set of real coding intents, a pattern corpus of
realistic size (200--2{,}000 artifacts) fails to route a large majority of
requests --- that is, if a substantial residue of intents shares no vocabulary
with any candidate artifact and resists paraphrase. I do not have this
experiment. The prototype exists to make it answerable, not to claim C1 is
established.

\section{Prior work and my relation to it}

This is not a new idea in the shape of the idea. It is a specific claim about
the \emph{unit} of reuse, and the literature contains most of the ingredients.

\textbf{Skill induction.} Work establishing that agents which memorise action
sequences fail, and that inducing \emph{rules} --- explicit control flow,
preconditions, variable binding --- generalises better (NSI, ICML 2026).
SkillGen synthesises a single auditable skill from trajectories. SkillRevise
couples execution evidence with general repair knowledge. My induction follows
this line: traces are grouped by the \textbf{primitive sequence}, and the
observed phrasings become aliases, so the induced unit is a rule and not a
script.

\textbf{Executable accumulation.} AgentFactory accumulates executable subagents
and reuses them. This is the closest prior work to my claim, and I claim no
novelty over it.

\textbf{Library learning and program synthesis.} DreamCoder-style systems induce
reusable libraries from traces. The difference is the target: a synthesised
program versus a declaration of intent bound to a pre-existing driver.

\textbf{Semantic caching.} Caches return prior \emph{responses}. Here the cache
entry is executable and the provider is not called. The relationship is a strict
generalisation.

My specific position is the combination of three properties: the reusable unit
is \textbf{content-addressed}, so two independently induced rules that are
textually identical collapse to one artifact without a comparison step; learning
is \textbf{zero-gradient}, producing inspectable artifacts rather than weight
deltas; and the objective is \textbf{escalation rate} rather than accuracy,
because accuracy is the wrong metric when a correct answer still costs a
request.

\section{The resource model}

Following \texttt{docs/ure-specification-formal.md}, a resource is a tuple
\[
  \mathcal{U} = \langle \mathcal{I}, \mathcal{S}, \mathcal{A}, \mathcal{C} \rangle
\]
where $\mathcal{I}$ is identity, $\mathcal{S}$ a state space, $\mathcal{A}$ a set
of action primitives, and $\mathcal{C}$ a set of constraints. Identity is a
content address: a DU-UUID derived from the canonical serialisation of the
manifest with \texttt{resource\_id} removed, hashed with SHA-256. Consequently a
resource cannot contain its own identity in its identity computation, and two
manifests that serialise identically \emph{are} the same resource.

\subsection{The gap between the formalism and the code}

The formal specification requires that constraints ``must return \texttt{True}
for an action to be dispatched to the Nucleus.'' The implementation prints
constraints for operator visibility and does not evaluate them. This is recorded
as divergence D5 and is not closed. I report it because it is load-bearing:
Section~\ref{sec:safety} assumes constraints gate dispatch, and today they do
not. The argument therefore rests on the evidence gate alone.

Two identity divergences also stand: the identifier scheme documented in prose
is not the one implemented (D2), and a document category vocabulary and a
resource-class enumeration are both in use for the same dimension (D8).

\section{Architecture}

\begin{center}
\small
\begin{tabular}{@{}ll@{}}
\toprule
\textbf{Tier} & \textbf{Behaviour} \\
\midrule
1 & exact / containment / IDF overlap; no model \\
1.5 & semantic re-rank of the residue; optional \\
2 & verified artifact executes; 0 tokens \\
3 & escalate to a provider; full cost \\
\bottomrule
\end{tabular}
\normalsize
\end{center}

The ladder's property is not accuracy but \textbf{monotone improvement}: every
promotion moves traffic from a paying tier to a free one, and the model is
retired by being out-routed rather than replaced. Tier 1 requires no model at
all, so the architecture is useful before any inference engine exists.

\subsection{Retrieval}

Tier 1 scores every candidate phrase against the intent using
inverse-document-frequency-weighted token overlap, with a decisive fast path for
containment. Trigram similarity is deliberately not used: it ignores term
frequency, and for short technical intents --- where one distinctive token
carries most of the signal --- that is strictly worse.

Measured, release build, single machine:

\begin{center}
\small
\begin{tabular}{@{}lrr@{}}
\toprule
\textbf{Corpus} & \textbf{p50} & \textbf{p99} \\
\midrule
12        & $4.6\,\mu$s   & $6.3\,\mu$s \\
1{,}000   & $405\,\mu$s  & $644\,\mu$s \\
5{,}000   & $2.17$ ms    & $3.85$ ms \\
20{,}000  & $9.96$ ms    & $11.2$ ms \\
\bottomrule
\end{tabular}
\normalsize
\end{center}

Scaling is linear at approximately $0.5\,\mu$s per artifact, which is the
limitation of a scan rather than of an inverted index. For comparison, SQLite
FTS5 with a real inverted index returns in $48\,\mu$s at 50{,}000 artifacts in an
18.6 MB file. Retrieval, not induction, is the component that breaks first at
scale.

\subsection{Storage as a view, not a copy}

The corpus is authoritative: a tree of \texttt{.ure} files under version
control. The analytical index is a \textbf{view over those files} rather than a
copy, so an edit in Git changes query results immediately and there is no
migration to write. Only mutable state --- promotions, traces, observations ---
is a table. This dissolves the consistency problem an ingest pipeline
introduces, at the cost of requiring the corpus to be schema-consistent, which
it is not yet.

\section{Induction: learning without gradients}

\texttt{src/induce} converts recorded interactions into candidate artifacts.
Traces are grouped by the \textbf{primitive sequence} the intent resolved to, not
by phrasing. Two sentences that reduced to the same primitives are evidence for
one rule; the sentences themselves become the action's aliases. A trace with no
recorded sequence is skipped rather than guessed at.

For each group the module synthesises a candidate with an action id equal to the
primitive sequence, stable across phrasings by construction; aliases drawn from
every distinct observed phrasing, plus both underscore and space forms;
constraints recording the observed failure rate, so a precondition is learned
rather than re-paid for on every call; and a \texttt{resource\_id} stamped with
the DU-UUID of the resulting body, so identical rules deduplicate by content
address.

It is not fine-tuning, and no component learns weights. Two consequences follow,
both intended: training requires no accelerator, being string and integer work
over a trace log; and the learned artifact is a readable file, so a wrong
artifact can be deleted and its deletion is a commit.

\section{Safety}
\label{sec:safety}

The characteristic risk of a system that promotes learned artifacts is that
\textbf{a wrong artifact becomes fast, silent and reused}. A slow wrong answer
is recoverable. A fast wrong answer served from a cache and re-derived from a
trace is not, because the trace then confirms it.

The prototype treats this as the design constraint rather than a downstream
check:

\begin{itemize}
\item Candidates are emitted as \texttt{candidate} and never as
      \texttt{champion}. Promotion is a separate, deliberate step; a function
      capable of emitting a promotion would have no gate at all.
\item A minimum of distinct observed phrasings is required. One observation is
      an anecdote, and the same sentence repeated ten times is one observation,
      not ten.
\item A maximum failure ratio blocks promotion. A sequence that usually fails is
      not a rule.
\item Failures become recorded preconditions rather than being discarded.
\item Traces without a recorded primitive sequence are not inducted from.
\end{itemize}

Confidence is multiplicative in breadth and reliability, and log-scaled in the
number of distinct phrasings so that it never saturates at the minimum. A
straightforward \texttt{observations / minimum}, clamped to 1, reads 1.0 the
instant the minimum is met, which would make the fourth distinct phrasing worth
exactly as much as the second; this was found by a test, not by inspection.

\subsection{A known weakness}

The evidence gate raises the cost of promotion; it does not establish
correctness. The evaluation reports a false-positive rate of 2 in 8 negatives,
and D5 means constraints do not yet gate dispatch. The honest summary is that
the system currently prevents \emph{thin} evidence from being promoted, and does
not yet prevent \emph{wrong} evidence from being promoted.

\section{Prototype evaluation}

A labelled fixture of 30 hand-written queries, split 20 expected-hit and 10
expected-negative. The corpus is the two curated manifests; the other ten are
quarantined and not indexed.

\begin{center}
\begin{tabular}{@{}lr@{}}
\toprule
\textbf{Metric} & \textbf{Value} \\
\midrule
Hit@1        & 18/20 --- \textbf{90.0\%} \\
Hit@3        & 18/20 --- 90.0\% \\
False negatives & 2 \\
False positives & 2 \\
Unsatisfiable  & 0 \\
\bottomrule
\end{tabular}
\end{center}

All four failures are manifest-authoring defects rather than semantic failures.
Three follow from aliases that are reasonable for their own action
(\texttt{open file}, \texttt{get content}, \texttt{create document}) and
actively harmful as resource-level match keys. The fourth requires either
constraint evaluation or semantic scoring.

One failure is not noise and is worth stating separately. IDF is currently
computed per \emph{phrase} across the corpus, so a resource listing
\texttt{valve} across many aliases has its own strongest term down-weighted for
appearing in its own alias set; a single rare term elsewhere then wins. The
corpus should be weighted per resource. This is unfixed, and it means the
reported 90\% is a floor rather than a current best.

\subsection{Cost against a local language model}

For the routing question specifically --- \emph{which artifact matches this
intent?} --- measured on the same machine:

\begin{center}
\small
\begin{tabular}{@{}lrr@{}}
\toprule
\textbf{System} & \textbf{Latency} & \textbf{Corpus} \\
\midrule
Qwen3-1.7B (128p+64o) & $\sim$2{,}860 ms & n/a \\
Tier 1 matcher & $\mathbf{4.6\,\mu}$s & 12 \\
Tier 1 matcher & $\mathbf{9.96}$ ms & 20{,}000 \\
Tier 1.5, bge-small & 8.3--20.8 ms & any \\
SQLite FTS5 & $48\,\mu$s & 50{,}000 \\
\bottomrule
\end{tabular}
\normalsize
\end{center}

The language model is roughly $6\times10^{5}$ slower than tier 1 at the current
corpus size. That comparison is not close, and it is also not the interesting
claim, because tier 1 does strictly less work. The interesting comparison is on
\emph{cost per served request}, where tiers 1 and 2 are free and tier 3 is not.

\subsection{Hardware}

All figures come from one machine: AMD Ryzen 7 5700U (Zen 2, 8c/16t), AMD Radeon
Vega 8 under Vulkan 1.4 (RADV RENOIR), 8 CUs, \textbf{no matrix cores}, sharing
15 GiB system RAM; reported by the backend as \texttt{matrix cores: none},
\texttt{int dot: 0}. This is a low-end mobile integrated GPU, reported in full
because throughput without hardware is not interpretable, and because the
absence of matrix cores is why decode throughput is roughly an order of
magnitude below published figures for the same model class on discrete GPUs.

Embedding tier, bge-small-en-v1.5 Q4\_K\_M, 33.21M parameters, 22.93 MiB:
967 t/s at 8 tokens, 3{,}032 t/s at 32, 6{,}158 t/s at 128. Generation tier,
Qwen3-1.7B: 272.8 t/s prompt processing and 26.8 t/s generation on the iGPU
against 155.1 and 22.8 t/s on CPU.

\section{Collection is distributed}

The first implementation of degradation kept every artifact's usage in a
supervisor-side table and ran a global planning pass. That is a
\emph{reporting} architecture, not an inference one, and it was wrong for this
system: \texttt{ActuatorDriver} exposes \texttt{execute()}, so an artifact is
itself the observation point. Routing that to a supervisor to obtain a decision
is strictly more work than deciding locally.

It also destroys a signal. An artifact attached to a nucleus that has gone away
is, to a central collector, indistinguishable from one that is genuinely idle.

The corrected design gives each artifact a self report --- connection, hit and
miss counts, distinct observations, last use, and reachability --- and a
decision function of exactly that report and a policy. No corpus argument, no
central usage table; the signature asserts the property.

\subsection{What distribution makes mandatory}

Moving the decision into the artifact exposed something the centralised version
had hidden: an artifact can have no usage because it is \textbf{unwanted}, or
because it is \textbf{unreachable}. Those require opposite responses.

The prototype already produces the second case. In the baseline,
\texttt{write output to results.txt} matched nothing, because the action id is
\texttt{write\_file} and containment requires the literal \texttt{write file}.
That pattern is not idle. It is mis-indexed, and it reports zero usage
indefinitely. A collector that cannot distinguish the two retires precisely the
artifacts that need repair, and does so silently, because the symptom of a
routing defect and the symptom of an unwanted artifact are the same observation.

Reachability is therefore a three-valued discriminator, and the third value is
the important one:

\begin{center}
\small
\begin{tabular}{@{}lll@{}}
\toprule
\textbf{State} & \textbf{Meaning} & \textbf{Action} \\
\midrule
Verified    & a query matched it & normal retirement \\
Unreachable & queried; no match & \textbf{quarantine} \\
Unknown     & never matched & \textbf{never retire} \\
\bottomrule
\end{tabular}
\normalsize
\end{center}

\texttt{Unknown} is the default, deliberately. Absence of a match is not
evidence of absence of demand, and most artifacts in a young corpus are in that
state; defaulting to \texttt{Unreachable} would empty the corpus within one
collection cycle.

Connection is authoritative for liveness and is not derived from lifecycle,
because the two are independent facts: a candidate that lost its nucleus is
served by lifecycle and detached by connection, and a harvested artifact a node
wired up is unverified but connected. Deriving one from the other made both
states unrepresentable and produced two sources of truth that could disagree.

Copies of one artifact disagree, because each observes only its own window. The
reconciliation rule is \textbf{monotone toward the less destructive action}: a
proposal to retire loses to evidence of use, always. A node with a narrow
window therefore cannot retire a capability a node with a broad window is
actively serving, and disagreement cannot compound. This is CRDT semantics
applied to artifact lifecycle, and the novelty claim is explicitly disclaimed.

Nothing about observation is centralised, but the evidence gate remains the
precondition for promotion. Distribution of observation does not require
distribution of trust: a mesh in which any attached node can synthesise and
promote arbitrary executables is precisely the attack surface described next.

\section{Exchange, redundancy, and safety}

\subsection{Structural deduplication, and its preconditions}

Exchange between models is a set union over content addresses, and is therefore
idempotent and commutative: two models that independently induce the same rule
exchange one artifact, not two. Redundancy is eliminated by the identity
function rather than detected after the fact. This holds only if three
properties are true, and one currently is not:

\begin{center}
\small
\begin{tabular}{@{}ll@{}}
\toprule
\textbf{Property} & \textbf{Status} \\
\midrule
Identical body $\to$ identical address & holds \\
Canonical serialisation byte-stable & by accident \\
Version label correct & \textbf{fails}: labelled v4 \\
\bottomrule
\end{tabular}
\normalsize
\end{center}

A v4 label is the worst available choice, because RFC~4122 defines v4 as
\emph{random}, and tooling that reasons about v4 semantics will mishandle these
identifiers. It should be v5, or a custom v8.

The canonicalisation property holds only by accident. The implementation comment
asserts that \texttt{serde\_json} does not guarantee key order; that comment is
wrong. Without the \texttt{preserve\_order} feature, \texttt{serde\_json::Map}
is \texttt{BTreeMap}-backed and the serialisation is already sorted. If any
contributor enables \texttt{preserve\_order} for an unrelated reason, the
comment becomes true and every content address in existence changes with no
error. This requires a test asserting address stability, not a comment.

\subsection{The corpus is an injection surface}

The characteristic risk of promoting learned artifacts is that a wrong artifact
becomes fast, silent and reused, and Section~\ref{sec:safety} addresses it from
the inside. Exchange adds a second path to the same failure, because a promoted
artifact carrying a payload \emph{is} a driver. A pattern that reaches champion
is persistent local code execution, and exchange is therefore a remote code
execution vector unless gated.

This is not hypothetical. koi Research reported ClawHavoc: 341 malicious skills
found by the bot they were targeting, rising to 824 by February 2026. unia's
exposure is worse than a typical skill library's, because
\texttt{ActuatorDriver::execute} is called by the Nucleus.

Four gates follow from structures already in the design: a
\texttt{payload\_kind: "none"} safety label, so the capability declaration and
the safety declaration become one field; re-verification of evidence on import,
never inheritance of a peer's; content addressing as integrity proof, so a
received manifest can be verified byte-identical to what the sender claims; and
mandatory lineage via \texttt{source\_external\_id}.

There is a further exposure the skill-library literature does not discuss. A
manifest's \texttt{guidance} and \texttt{aliases} are free text. If any part of
a manifest is ever placed in a model's context --- which is what a pattern-aware
agent would do --- then an attacker who controls the corpus controls part of the
prompt. That is prompt injection arriving through the retrieval store rather
than through the user.

\subsection{Union is not safe when precision binds}

The natural design is to merge pattern stores as a set union. The evaluation
shows why that is wrong for this system. Precision is the binding constraint at
2 false positives across 10 negatives, and precision \emph{degrades} as the
corpus grows, because every alias is an additional opportunity to match the
wrong artifact. Accepting a peer's patterns adds aliases faster than it adds
coverage.

So a model can get \textbf{worse} by receiving patterns, and union is harmful
exactly when it appears most attractive. Exchange must be evidence-weighted:
artifacts are exchanged \emph{with} their observation count, failure ratio and
provenance; the evidence gate is re-run on import rather than trusted; and
provenance is mandatory. This is a negative result, and the most useful claim in
this section, because it contradicts the obvious design.

Three further mechanisms are needed beyond content addressing. Aliases should
be edges in a graph, \texttt{alias} $\rightarrow$ (artifact, action), rather than
fields owned by an artifact: two artifacts can then share an alias without
duplicating it, and the alias is scoped to an action rather than a resource,
which is the direct cause of two of the four failures above. Patterns should
compose over other patterns rather than restate their steps. And promotion
should be monotonic --- if a champion already covers a capability, induction
refines it rather than creating a sibling, or every re-induction cycle produces
near-duplicates differing only in aliases, which never dedupe.

\subsection{Interoperability, by severity}

The constraint grammar is undefined: \texttt{constraints} is an array of free
strings, two implementations will parse it differently, and D5 means neither
evaluates it. You cannot execute a constraint you cannot parse, so this blocks
cross-model execution outright. It is a specification task, not an engineering
one.

The remaining gaps are degraded mode rather than blockers: no capability
declaration, so a receiver cannot be told which of the 20 primitives it
implements; no schema negotiation, since \texttt{ure\_version} is decorative and
nothing reads it; the DU-UUID version mislabelling above; unspecified canonical
serialisation, with one code comment wrong about it; ambiguous alias semantics,
resource-level or action-level, which is the scoring bug; and no manifest JSON
Schema, so nothing validates an import.

\section{What is not demonstrated}

I consider this section load-bearing rather than a disclaimer.

\begin{itemize}
\item \textbf{No escalation rate is measured.} The corpus holds two usable
      artifacts. A rate over two artifacts is not an escalation rate.
\item \textbf{The fixture is 30 queries authored by the system's author.} It is
      a smoke test. It found a real defect in an afternoon, which is its
      purpose, but it cannot estimate behaviour on real traffic. A 200-query
      fixture written by someone who did not build the matcher is the minimum
      credible next step.
\item \textbf{No learned component has run.} The induction module is
      implemented and covered by tests, but no trace log is being written:
      \texttt{unia\_record} exists and no agent calls it.
\item \textbf{Collection has never run.} The policy is tested against
      synthetic usage and has never been reconciled against a real
      disagreement, so the monotonicity argument is unvalidated.
\item \textbf{Constraints do not gate dispatch} (D5).
\item \textbf{The analytical index is non-functional.} \texttt{term\_idf} in
      the DuckDB retrieval layer raises a type-coercion error; the corpus views
      load and are verified, the ranking path is not.
\item \textbf{The corpus is not schema-consistent.} Twelve artifacts carry
      roughly a dozen distinct key sets, one retains a corrupted key from an
      earlier substitution defect, none carries a human-readable title, and two
      use human-assigned identifiers where the specification expects a UUID.
\item \textbf{No comparison against a learned router} exists, because no
      learned router exists.
\item \textbf{The four failures above are unfixed}, so the reported 90\% is a
      floor.
\end{itemize}

\section{Threats to validity}

The author wrote both the fixture and the matcher, which is a smoke test that
found a real defect but is not an unbiased evaluation. $n = 2$ artifacts, so
everything measured is at a corpus size with no statistical meaning. The safety
argument has one pillar, since D5 leaves constraints ungated. No learned
component has run against real usage, and its behaviour on adversarial or
degenerate traces is unknown. The prior-art claims name the closest work I am
aware of; a reviewer with better recall may find more, and I would expect a
non-trivial amount of it. The prototype is 45 tests and one fixture, which is
evidence of internal consistency, not external validity. Every figure comes
from a Vega 8 with no matrix cores, the least favourable device in the class;
the relative claims should transfer, the absolute latencies will not.

\section{Contributions}

Stated conservatively, relative to the prior work surveyed above.

\begin{enumerate}
\item A formalisation of the reusable unit as an executable, content-addressed
      artifact, with the property that identical induced rules deduplicate
      without comparison and renaming does not change identity.
\item A zero-gradient induction mechanism deriving both rule and alias set from
      recorded interactions, with an evidence gate specified to resist promotion
      on thin evidence.
\item Escalation rate as the primary objective, with the argument that accuracy
      is the wrong metric when a correct answer still costs a request.
\item A three-valued reachability discriminator arising from the observation that
      collection must be distributed, and a retirement protocol that is monotone
      toward the less destructive action.
\item A negative result on exchange: content-addressed union is unsafe when
      precision binds, so a model can degrade by receiving patterns.
\item A measured characterisation of the lexical and semantic tiers on low-end
      integrated hardware, including the observation that a scan is adequate to
      $\sim 10^{4}$ artifacts and an inverted index is required beyond it.
\item A first retrieval-quality baseline locating current failure modes in
      manifest authoring rather than retrieval, and identifying a
      document-frequency weighting defect that is cheap to fix.
\end{enumerate}

I do not claim novelty over executable skill accumulation, a working learned
model, a measured escalation rate, novelty in applying CRDT semantics to
artifact lifecycle, or that the safety argument is complete.

\section{Future work, in dependency order}

\begin{enumerate}
\item Fix the manifest defects surfaced by the evaluation --- scope aliases to
      actions, add the space forms, and correct document-frequency weighting to
      the resource level. Expect this to move the fixture to 100\% and to
      establish how much of the remaining problem is authoring rather than
      retrieval.
\item Close the induction loop: write real traces, induce, re-measure. This is
      the first end-to-end exercise of the learning loop and the first test of
      whether induced aliases beat authored ones.
\item Falsify or support C1 with a fixture of 200 or more intents authored
      independently, reporting escalation rate and false-positive rate.
\item Define the constraint grammar, so constraints gate dispatch and the safety
      argument gains its intended second pillar.
\item Only then evaluate a semantic tier, and only on the residue that survives
      the first three steps. If the residue is small, the correct decision is to
      stop.
\item Define what a correct pattern is. Every step above assumes an oracle for
      correctness; self-evaluation without a verifier is the unsolved part, and
      property-based testing against observed outcomes is the current approach.
\end{enumerate}

\section*{Authorship and assistance}

This paper is a single-author work by \textit{iulian dacineu}, authored and
researched with AI assistance. The author line uses a pseudonym. An earlier
project by the same author, \emph{eunia} --- European Union National Informative
Assistant --- was submitted to the European Commission's EUvsVirus Hackathon on
24 April 2020, and is cited by its canonical URL; the present work is a distinct
design and not a continuation of it, although the distributed and
multidimensional framing originates there. The prototype is MIT licensed;
commercial enquiries are directed through \texttt{LICENSE-COMMERCIAL.md}.

\begin{thebibliography}{9}

\bibitem{nsi} Lifting Traces to Logic: Programmatic Skill Induction with
  Neuro-Symbolic Learning for Long-Horizon Agentic Tasks. ICML 2026.

\bibitem{skillgen} SkillGen: synthesis of auditable skills from trajectories via
  contrastive induction.

\bibitem{skillrevise} SkillRevise: trace-conditioned skill revision coupling
  execution evidence with general repair knowledge.

\bibitem{agentfactory} AgentFactory: a self-evolving framework through
  executable subagent accumulation and reuse.

\bibitem{sok} SoK: Agentic Skills --- Beyond Tool Use in LLM Agents, 2026.

\bibitem{dontbreak} Don't Break the Cache, arXiv:2601.06007, 2026.

\bibitem{bge} BAAI. \emph{bge-small-en-v1.5} model card.

\bibitem{rfc4122} RFC 4122. \emph{Universally Unique IDentifiers (UUIDs)}.

\bibitem{spec} \emph{unia} specification and divergences:
  \texttt{docs/SPEC.md}, \texttt{docs/ure-specification-formal.md}.

\end{thebibliography}

\end{document}
